\section{Commuting fields and scalar-twisted deleted Heisenberg carriers}
\label{sec:heisenberg-commuting-field-carrier}

\paragraph{Author freeze and provenance.}
D-TD34--36, reserved question Q2er, strictly extends the qualified
D-TD24--26 argument to defining modules which need not be absolutely
irreducible over the prime field.  Root proposed this extension;
\texttt{/root/g34\_eightfolds} supplied the uniform commuting-field
bound, the scalar-twisted deleted-module argument and the full proof.
This separate section does not alter the earlier frozen proof.
The complete author proof, source-interface, checker and first-output
readings pass.  Independent and root whole-proof reviews also pass.

\paragraph{Actual setup.}
Let an actual connected finite $J$-cover of the sphere have a finite
coefficient representation on $V$, with a stable whole factorization
$V=W\otimes R$ and the original inertia budget
\[
 e=\dim V,\qquad
 \sum_i\operatorname{rank}(1-V(g_i))=2e+6.
 \tag{HCF.1}
\]
The entire external partner $R$ is retained.  Suppose the actual
projective image $P=\operatorname{im}(J\to\operatorname{PGL}(W))$
contains a normal nondegenerate irreducible projective Heisenberg label group
$Q\simeq\mathbf F_p^{2m}$, and write
\[
 H=P/Q\hookrightarrow\operatorname{Sp}(Q),\qquad
 \psi:H\twoheadrightarrow T\in\{A_n,S_n\},\qquad n\ge10^6.
 \tag{HCF.2}
\]
The defining $\mathbf F_p H$-module $Q$ is assumed irreducible.
Every map here is an actual map of the indicated finite images.
A subgroup or composition factor is not substituted for $\psi$.

Schur's lemma and the finite division-ring theorem make
$F=\operatorname{End}_{\mathbf F_pH}(Q)$ a finite field.
Put
\[
 F=\mathbf F_q,\qquad q=p^f,\qquad
 a=\dim_F Q,\qquad af=2m,\qquad \dim_{\mathbf C}W=p^m=q^{a/2}.
 \tag{HCF.3}
\]
Commutativity removes the opposite-ring convention for the action of
$F$.  Every element of $H$ is $F$-linear.  Moreover $Q$ is absolutely
irreducible in this $a$-dimensional $F$-model.  Indeed, the finite image
algebra $\mathbf F_p[H]|_Q$ acts faithfully and irreducibly, has zero
nilpotent radical, and is a single full matrix algebra over a finite
division ring.  Its simple module and commuting algebra identify that
field with $F$ and the image with $\operatorname{Mat}_a(F)$.
Scalar extension of this full matrix algebra proves absolute
irreducibility.  These are precisely the finite-algebra statements
and proofs in Knapp, Chapter~II, Theorems~2.2, 2.13 and~2.48,
pp.~83--84, 91--92 and 117--118.\cite{TD:Knapp2016}
No primitivity or tensor indecomposability follows from this argument.

\subsection*{D-TD34. A uniform bound on the intrinsic commuting field}

\begin{proposition}[D-TD34]
In (HCF.1)--(HCF.3), without a primitivity or tensor hypothesis,
\[
                  a\ge2,\qquad q<28900,
                  \qquad f\le14.
 \tag{HCF.4}
\]
\end{proposition}

\begin{proof}
If $a=1$, the actual $H$ embeds in $F^*$ and is abelian, contrary
to (HCF.2).  Every $F$-scalar in $H$ is central in $H$ and maps
to the trivial center of $T$.  Thus $\psi(h)\ne1$ implies that
the $F$-matrix of $h$ is nonscalar, and in particular
\[
 \operatorname{rank}_{\mathbf F_p}(1-h)
       =f\operatorname{rank}_{F}(1-h)\ge f.
 \tag{HCF.5}
\]
The accepted exact D-TD05 Weyl-basis identity therefore gives
normalized absolute trace at most $c=q^{-1/2}$ for every unitary
Heisenberg normalizer lift whose $T$-image is nonidentity.  The
entire external partner cannot increase that normalized trace.

We retain all original lift periods.  For $g\in J$, let
$L=\operatorname{ord}_J(g)$ and let $o$ be the order of its
image in $T$.  Then $o\mid L$, and among the powers numbered
$0,\ldots,L-1$, exactly $L/o$ have trivial $T$-image.  Bound
their normalized absolute traces by one, and every other power
by $c$.  Fourier inversion with period $L$ gives, for every
target eigenvalue and hence for actual moving rank,
\[
 \frac{\operatorname{rank}(1-V(g))}{e}
                \ge(1-c)(1-1/o).
 \tag{HCF.6}
\]
This is valid whether or not the coefficient operator itself
has order $L$.  It does not assert that $W$ descends to $T$.

The actual connected $T$-quotient of the sphere cover has genus
at least two, since $T$ is nonsolvable and has order greater
than sixty.  The complete genus-zero/genus-one exclusions and
positive orbifold enumeration of D-TD02 give
$\sum_i(1-1/o_i)\ge85/42$, with identity images omitted.
Summing (HCF.6) yields
\[
                  e(1-85q^{-1/2})\le252.
 \tag{HCF.7}
\]
If $q\ge28900=170^2$, this forces $e\le504$.  But (HCF.3)
and $a\ge2$ give $e\ge q^{a/2}\ge q>504$, a contradiction.
Thus $q<28900$.  Since $2^{15}>28900$, also $f\le14$.
The field bound uses neither the absolute Guralnick--Saxl
hypotheses nor a lower bound on the number of projected marks.
\end{proof}

\subsection*{D-TD35. Deleted modules with arbitrary scalar twists}

For a field $F$ of characteristic $p$, write
\[
 D_n(F)=\{(v_i)\in F^n:\sum_i v_i=0\}
       \big/\bigl(F(1,\ldots,1)\cap\{\sum_i v_i=0\}\bigr).
 \tag{HCF.8}
\]
It has dimension $n-1$ if $p\nmid n$ and $n-2$ otherwise.
Its ordinary natural permutation operator is denoted by $D_y$.

\begin{lemma}[D-TD35]
Suppose, in the actual setup above, that an $F$-linear identification
$Q\simeq D_n(F)$ has the property
\[
                    h=\lambda(h)D_{\psi(h)},\qquad
                    \lambda(h)\in F^*,\qquad h\in H.
 \tag{HCF.9}
\]
Then the sphere budget (HCF.1) is impossible.
No restriction to characteristic two, trivial scalar kernel, or
particular invariant $F$-bilinear form is imposed.
\end{lemma}

\begin{proof}
For any permutation $y$ and nonzero scalar $\lambda$, the kernel
of $1-\lambda y$ on $F^n$ has dimension at most the number of
cycles of $y$: each cycle contributes either zero or one free
coordinate, according as its consistency equation holds.
This is valid in defining characteristic, without semisimplicity.
Restriction to the augmentation hyperplane loses at most one
from the rank, and its further quotient in (HCF.8) loses at
most one more.  Therefore
\[
 \operatorname{rank}_{F}(1-\lambda D_y)
                  \ge\operatorname{ind}_n(y)-2.
 \tag{HCF.10}
\]
For $y\ne1$ that rank is also at least one.  Here is an elementary
projective-faithfulness check.  If $D_y$ were scalar, apply it to
every coordinate difference $e_i-e_j$.  The difference between
the resulting vectors in the augmentation space is constant
modulo the line in (HCF.8), but has support at most four.
Since $n\ge5$, that constant is zero.  Thus $y$ preserves every
unordered pair of letters and is the identity.  The same argument
works in characteristic two.  It follows that (HCF.9) is nonscalar
whenever $\psi(h)\ne1$.

Put $k=n-\operatorname{fix}_n(y)$.  Multiplying the rank in
(HCF.10) by $f$ and applying the exact Heisenberg trace identity
gives the following bounds for every lift above $h$.
If $y\ne1$ and $k\le n/16$, its normalized absolute trace is
at most
\[
 q^{-1/2}\le2^{-1/2}< (15/16)^3
                  \le\bigl(\operatorname{fix}_n(y)/n\bigr)^3.
 \tag{HCF.11}
\]
The strict middle inequality follows by squaring:
$2\cdot15^6>16^6$.
If $k>n/16$, then $\operatorname{ind}_n(y)\ge k/2$ and
\[
 \frac{|\operatorname{tr}A|}{\dim W}
    \le p^{-\operatorname{rank}_{\mathbf F_p}(1-h)/2}
    <2^{1-n/64}<n^{-3}.
 \tag{HCF.12}
\]
For the last inequality we use the elementary dyadic estimate
proved in D-TD22, valid already for every $n\ge4096$.
Indeed $\operatorname{rank}_{\mathbf F_p}(1-h)
\ge f(\operatorname{ind}_n(y)-2)>n/32-2$ in this range.
Together with the trivial estimate at $y=1$, these prove the
full cubic bound using $\max\{\operatorname{fix}_n(y),1\}$.
On every nonidentity natural image the uniform trace bound
is at most $1/\sqrt2<7/8$.

All original powers and reciprocal tensor scalars remain present.
The finite projective factor image has a finite special-linear
lift, and its representatives can be taken unitary.  For an
original operator $V(g)$, choose such a representative $A$ on
$W$ and put the reciprocal scalar into $B$ on $R$, so that
$V(g)=A\otimes B$.  If $L=\operatorname{ord}_J(g)$, then
$A^L=\zeta I$ for a root of unity and $B^L=\zeta^{-1}I$.
Thus $B$ is finite order, and
$|\operatorname{tr}B^j|/\dim R\le1$ for every integer $j$.
The bounds above apply to each $A^j$ and the actual natural
image $\psi(h)^j$.  Consequently the full coefficient character
satisfies exactly the full-power hypotheses of HD-FT27.

Finally $\dim_F Q\ge n-2$, so
\[
 e\ge\dim W=q^{\dim_F Q/2}
         \ge2^{(n-2)/2}>n^3.
 \tag{HCF.13}
\]
The last elementary inequality is also part of D-TD22 at
$n\ge4096$.  To match the faithful-image formulation of HD-FT27,
put $K=\ker V$.  An element of $K$ acts projectively trivially
on $W$, since a tensor operator equal to the identity has scalar
factors.  Thus $K$ kills $P$ and $T$.  The actual connected
quotient cover $C/K\to\mathbf P^1$ has group $J/K$, on which
$V$ is faithful.  Removing identity coefficient inertia images
does not change their moving-rank sum.  All character bounds
descend, as does the actual surjection to $T$.
Apply accepted HD-FT27 to this actual faithful coefficient image.
It excludes (HCF.1) at $n\ge10^6$ with any entire partner.
The actual alternating form on the original $\mathbf F_p$-space
is retained throughout; no identification with the dot product
as an $F$-bilinear alternating form was used.
\end{proof}

\subsection*{D-TD36. The qualified commuting-field exclusion}

\begin{theorem}[D-TD36]
In (HCF.1)--(HCF.3), assume that the $a$-dimensional $F$-model
of $H$, after scalar extension to an algebraic closure of $F$,
is primitive and tensor-indecomposable in the precise
Guralnick--Saxl sense.  Then there is no such actual sphere source.
The original $\mathbf F_p$-module need not be absolutely irreducible.
Here tensor indecomposability excludes a factorwise preserved tensor
product; a permutation of equal tensor factors is not silently excluded.
\end{theorem}

\begin{proof}
The $F$-model is already absolutely irreducible by the finite
algebra argument following (HCF.3).  D-TD34 gives $q<28900$.
If $a<576$, then $a\le575$, and the actual faithful matrix
group satisfies
\[
 |H|\le|\operatorname{GL}_a(q)|<q^{a^2}
       <2^{15\cdot575^2}=2^{4959375}
       <2^{8999999}<\frac{n!}{2}\le|T|.
 \tag{HCF.14}
\]
As in (GHC.9), the last $500000$ factors of $1000000!$ are
each greater than $2^{18}$, proving the last strict comparison.
This contradicts surjectivity, so $a\ge576$.

Apply Guralnick--Saxl Theorem~7.1 and Remark~3 to this
$a$-dimensional representation over the algebraic closure, not
to the $af$-dimensional scalar restriction.\cite{BI:GS}
In its generic branch every nonscalar $h$ has $F$-rank at
least thirteen.  Quotient-nonidentity elements are nonscalar,
so their Heisenberg trace ratios are at most
$2^{-13/2}<1/90$.  The same full-period/actual-$T$-cover proof
as (HCF.6)--(HCF.7) then gives $e<4536$, whereas
$e\ge q^{a/2}\ge2^{288}>4536$.  Hence this branch is impossible.

Since $a>10$, the only remaining source exceptions are the natural
classical and deleted alternating/symmetric cases.  We give the
exact normal-component interface.  The supplied source proof first
identifies $S=F^*(H/Z(H))$ as simple in these cases.  The stable
derived subgroup $E$ of its central preimage is normal and
quasisimple.  Over the algebraic closure, the $E$-socle is
nonzero and $H$-invariant, hence is the whole space.  Primitivity
makes its restriction homogeneous.  Its irreducible $E$-factor
has dimension greater than one, and the evaluation tensor
decomposition has both factor algebras normalized by $H$.
Tensor indecomposability makes its multiplicity one.  Thus $E$
acts absolutely irreducibly on the $F$-model.

Its centralizer in $H$ is the scalar center.  An automorphism
of a perfect central extension inducing the identity on its
simple quotient is trivial, since the difference is a central
homomorphism.  Therefore $H/(EZ(H))$ embeds in
$\operatorname{Out}(S)$, and $H/E$ is solvable by the exact
published Schreier input.\cite[Theorem2.3, p.~1314]{HD:BabaiLuksSeress1997}
The perfect normal image $\psi(E)$ cannot be trivial and so
equals $A_n$.  Thus $S\simeq A_n$.
The complete classical/alternating repetition list in Li--Xia,
Section~2.2.1, numbered author p.~13, contains only $A_5,A_6,A_8$;
it excludes the full natural classical core here.\cite{HES:LiXia2022}
This argument uses the actual full normal core, not an arbitrary
alternating subgroup of a classical group.

In the deleted case $E$ is the actual ordinary $A_c$ on its
fully deleted module.  The derived group of a scalar extension
of this ordinary perfect group is that group itself, so scalar
choices introduce no new central cover.  Its simple image
$A_n$ gives $c=n$.  In particular $\psi|_E$ is injective.
For $h\in\ker\psi$, its commutators with $E$ belong both
to $E$ and to $\ker\psi$, hence are trivial.  Such an $h$
is therefore an $F$-scalar.  We retain this possible central
cyclic kernel; it need not vanish, and the characteristic
need not be two.

The identification of $E$ through $\psi$ differs from the
standard alternating identification by an automorphism of $A_n$.
Since $n\ge10^6$, that automorphism is induced by a permutation
of its letters.  After relabeling, the $E$-module is isomorphic
over the algebraic closure to $D_n(F)$, which is defined already
over the prime field.  The latter is absolutely irreducible:
the coordinate-difference proof in D-TD25 works in every
characteristic and in this same natural degree.

The intertwiner equations are linear over $F$.  Their solution
space has dimension one after scalar extension, hence already
has a nonzero $F$-solution.  It is an isomorphism.  Under that
identification, each actual $h\in H$ and $D_{\psi(h)}$
implement the same conjugation on the absolutely irreducible
core.  Their ratio is an $F$-scalar.  We have obtained precisely
(HCF.9) for every element, with the actual quotient $\psi$.
Multiplication even makes $\lambda:H\to F^*$ a character,
although only its scalar values and their powers are needed.
No splitting of the possible central extension is assumed.
D-TD35 now gives the contradiction.
\end{proof}

\paragraph{The homogeneous irreducible-kernel branch.}
If $N=\ker(H\to A_n)$ is irreducible on the original $Q$,
then $F_N=\operatorname{End}_N(Q)$ is a finite field.
Its conjugation action by $H$ factors through the perfect
$A_n$ and a cyclic finite-field automorphism group, so is
trivial.  Thus $F_N\subseteq\operatorname{End}_H(Q)$;
the reverse inclusion is automatic and they are equal.
D-TD36 therefore applies to this branch whenever the stated
absolute primitivity and tensor indecomposability of its
commuting-field model hold.  Neither condition is inferred
from irreducibility or from its prime-field analogue.

\paragraph{Strictness of the field extension.}
The larger module domain is nonempty even in the stated large
natural degrees.  Here is a defining-group example, not an
example satisfying the forbidden curve budget.  Take $p=3$,
$F=\mathbf F_9$, $n\ge10^6$ with $3\nmid n$, and
$D=D_n(F)$, so $a=n-1$.  Let $Z\simeq C_4$ be the norm-one
scalars in $F^*$, and let $H=Z\times A_n$ act by scalar
multiplication and the ordinary deleted representation.
This action is faithful: the deleted alternating action is faithful,
and its scalar intersection is trivial because $A_n$ has trivial center.
Its generated $\mathbf F_3$-algebra is
$\operatorname{Mat}_{n-1}(\mathbf F_9)$: the alternating core
generates the full prime-field matrix algebra, and a generator
of $Z$ generates $\mathbf F_9$ over $\mathbf F_3$.
Thus the original representation is irreducible with commuting
field $F$, and is not absolutely irreducible over $\mathbf F_3$.

Choose $0\ne\kappa\in F$ with $\bar\kappa=-\kappa$, where
the bar is the nontrivial field automorphism.  The form
\[
       (v,w)\longmapsto
       \operatorname{Tr}_{F/\mathbf F_3}
              \left(\kappa\sum_i v_i\bar w_i\right)
 \tag{HCF.15}
\]
on the augmentation space is nondegenerate and alternating
over $\mathbf F_3$, and is preserved by $Z\times A_n$.
Nondegeneracy uses $3\nmid n$, and alternation follows because
$\sum_i v_i\bar v_i$ is fixed by the bar and
$\operatorname{Tr}(\kappa)=0$.

The $F$-model is absolutely primitive: a nontrivial imprimitivity
system has at most $n-1$ blocks, so the simple $A_n$ acts
trivially on them by factorial orders, contradicting its
irreducibility.  It is absolutely tensor-indecomposable as well.
If both factors of a nontrivial preserved tensor product have
dimensions $u,v$, the two projective $A_n$-actions must be
nontrivial, or core irreducibility fails.  Each is faithful
and has degree greater than $n/5$ by D-TD32.  Thus
$uv>n^2/25>n-1=a$, a contradiction.  This verifies a strict
enlargement of the defining-module hypotheses beyond D-TD24--26.
It constructs no actual source or successful polarization.

\paragraph{Retained boundary.}
Absolute primitivity and tensor indecomposability are required
for the representation of dimension $a$ over its commuting field,
after scalar extension.  Imprimitive and tensor-decomposable
absolute constituents remain outside D-TD36.  General $S_n$
sign actions on kernel endomorphism fields are not identified
with the homogeneous $A_n$ case.  No bound on all remaining
Heisenberg half-ranks, no independent lift to an original
characteristic-zero group, and no equality of original character
fields or polarizations is asserted.
